\scanpage{58}
\mspage{51}

$W$ est \uline{quillénisable sur $\mathrm{Diag}$}, si pour tout $I$ dans $\mathrm{Diag}$, \deleted{\unsure{le}} $(M(I),W(I))$ est quillénisable.

% slanted

\margin{Je supposerai que $\mathrm{Diag}$ contient \deleted{\unsure{tout}} au moins les cat. provenant d'ens\unsure{.} ordonnés finis.}

En pratique, j'utiliserai cette \unsure{locution|notion} dans le cas où $\mathrm{Diag}=\mathrm{Cat}$, ou alors, qu'il \unsure{satisfasse} au moins les propriétés de stabilité utiles pour le domaine d'un dérivateur : stable par \deleted{\unsure{prod}}\,$\varprojlim$ finies, par sommes finies -- éventuellement par $\underset{\unsure{-|\to}}{\lim}$ finies, etc.

\uline{Exemple \uline{1}} : Celui des § 3, 4, où on a dégagé des conditions sur $(M,W)$, avec $M$ une ``\unsure{grosse|pure}'' catégorie et $M$ \unsure{accessible}, pour que \unsure{lui|$W$} soit quillénisable sur $\mathrm{Cat}$ tout \unsure{entier|entière}. C'est le cas qui me paraît à présent le plus intéressant.

\uline{Exemple 2} : Soit\deleted{\unsure{ent}} $A$ une catégorie abélienne, et $M=K^{\unsure{\bullet|\cdot|\prime}}(A)$ la catégorie des complexes de \unsure{cochaines|cochaînes} à valeurs dans $A$, $W$\inserted{$=W_A$} l'ens. des quasi-isomorphismes, donc

\[ MW^{-1}=K^\bullet(A)\,W_A^{-1}\simeq D(A)\quad\text{catégorie dérivée de }A. \]

On définit \unsure{aussi} $M^+,M^-,M^b$, avec les localiseurs induits $W_A^+,W_A^-,W_A^b$, et les cat. de fractions $D^+(A),D^-(A),D^b(A)$. On
% archive sheet number: 57
